% Chương 1 của đề (cài đặt). Nguồn: docs/report/chapter1.md; số liệu ở results/metrics/optimization_*.
\section{Cài đặt Eclat}
\label{sec:impl}

\subsection{Định nghĩa và cơ sở thuật toán}

Cho cơ sở dữ liệu (CSDL) $D$ gồm $N$ giao dịch, mỗi giao dịch là một tập item. Với itemset $X$, gọi $T(X)$ là
TID-list của $X$, tức tập chỉ số (TID) của các giao dịch chứa $X$. Support tuyệt đối của $X$ là
\begin{equation}
\sigma(X) = |T(X)|.
\label{eq:support}
\end{equation}
Với ngưỡng tương đối $s \in (0, 1]$, cài đặt dùng ngưỡng tuyệt đối
\begin{equation}
\theta = \lceil s N \rceil,
\label{eq:threshold}
\end{equation}
và xuất mọi itemset $X \neq \emptyset$ thỏa $\sigma(X) \geq \theta$ cùng support của nó. Cách làm tròn lên
khớp với SPMF: trên CSDL mẫu 5 giao dịch, $s = 21\%$ cho $sN = 1{,}05$ và $\theta = 2$, ra 15 itemset như SPMF;
làm tròn thông thường sẽ cho $\theta = 1$ và 19 itemset.

Tính chất nền tảng của Eclat~\cite{zaki2000scalable} là
\begin{equation}
T(X \cup Y) = T(X) \cap T(Y).
\label{eq:intersect}
\end{equation}
Phép giao không làm TID-list dài thêm, nên mọi itemset chứa một itemset không phổ biến cũng không phổ biến; thuật toán
dùng tính chất này để tỉa nhánh.

\subsection{Bản cơ bản với TID-list}

Dữ liệu ngang có kiểu \texttt{Vector\{Vector\{Int\}\}}, mỗi giao dịch là các item nguyên không âm tăng
nghiêm ngặt. Bản cơ bản (\texttt{src/algorithm/baseline.jl}) chuyển dữ liệu sang ánh xạ từ item tới
\texttt{Vector\{Int\}} chứa các TID tăng dần, loại item đơn không đủ support, sắp các item còn lại tăng dần
rồi duyệt sâu theo tiền tố (Thuật toán~\ref{alg:dfs}). Một lớp với tiền tố $P$ là danh sách $E$ các cặp
$(x, T_x)$, trong đó $x$ là item mở rộng và $T_x = T(P \cup \{x\})$; $C$ là lớp con được dựng cho tiền tố
$P \cup \{x\}$. Lời gọi đầu tiên có $P = \emptyset$ và $E$ gồm các item đơn phổ biến đã sắp.

\begin{algorithm}[t]
\caption{Mở rộng lớp tiền tố}
\label{alg:dfs}
\begin{algorithmic}[1]
\Procedure{DFS}{$P$, $E$}
  \For{$(x, T_x)$ tại vị trí $i$ trong $E$}
    \State thêm $x$ vào $P$; xuất $P$ với support $|T_x|$
    \State $C \gets \emptyset$
    \For{$(y, T_y)$ đứng sau vị trí $i$}
      \State $T \gets T_x \cap T_y$
      \If{$|T| \geq \theta$} thêm $(y, T)$ vào $C$ \EndIf
    \EndFor
    \If{$C \neq \emptyset$} \Call{DFS}{$P$, $C$} \EndIf
    \State bỏ $x$ khỏi $P$
  \EndFor
\EndProcedure
\end{algorithmic}
\end{algorithm}

Phép giao hai danh sách đã sắp dùng hai con trỏ: con trỏ ở giá trị nhỏ hơn tiến lên; khi hai giá trị bằng nhau
thì ghi nhận TID một lần và tiến cả hai. Ví dụ với $D = [\{1,2\}, \{1\}, \{1,2\}]$, ta có $T(1) = [1,2,3]$,
$T(2) = [1,3]$ và $T(\{1,2\}) = [1,3]$; với $\theta = 2$, kết quả là $\{1\}{:}3$, $\{1,2\}{:}2$, $\{2\}{:}2$.

Mỗi tổ hợp chỉ được mở rộng theo thứ tự item tăng dần nên không sinh trùng. Mọi tập con của một itemset phổ
biến đều phổ biến, nên các tiền tố cần thiết không bị tỉa mất; vì vậy thuật toán sinh đủ và đúng các itemset
phổ biến.

\subsection{Bản tối ưu với bitset}

Bản tối ưu (\texttt{src/algorithm/optimized.jl}) giữ cùng cách phân lớp và duyệt sâu, nhưng thay TID-list bằng
\texttt{BitVector} dài $N$ bit: bit thứ $t$ bằng 1 khi giao dịch $t$ chứa itemset. Phép giao
(\ref{eq:intersect}) thành phép AND từng bit, và support là số bit 1. Chỉ item đơn đủ ngưỡng mới được cấp mask;
mỗi phép AND tạo một mask mới, không sửa mask cha hay mask anh em. Một từ máy 64 bit chứa 64 TID nên phép giao
rẻ hơn giao TID-list khi support cao; ngược lại, với dữ liệu thưa, việc cấp và quét mask dài $N$ cho cả những phép giao thất bại
có thể tốn bộ nhớ hơn TID-list ngắn. Bitset không luôn tốt hơn ở mọi tiêu chí.

\subsection{Giao diện, nhập xuất và dòng lệnh}

Hai hàm \texttt{eclat\_baseline} và \texttt{eclat\_bitset} cùng nhận CSDL và ngưỡng tuyệt đối, trả về danh sách
cặp (itemset, support); hàm \texttt{eclat} mặc định gọi bản bitset. Hàm đọc tệp theo định dạng SPMF: tách item
bằng khoảng trắng, sắp và bỏ item trùng trong từng dòng, giữ các giao dịch lặp, bỏ dòng trống và dòng ghi chú
bắt đầu bằng \texttt{@}, \texttt{\#}, \texttt{\%}; dòng chỉ chứa khoảng trắng được giữ như một giao dịch rỗng,
giống SPMF. Hàm ghi xuất mỗi dòng dạng \texttt{i1 i2 ... \#SUP: n} với support tuyệt đối.

Ngưỡng $s$ được đọc thành số hữu tỉ (ví dụ chuỗi \texttt{"7\%"} thành $7/100$) thay vì số thực, vì
$0{,}07 \times 100$ trong số thực bằng $7{,}000000000000001$ và (\ref{eq:threshold}) sẽ cho 8 thay vì 7.
Chương trình dòng lệnh chỉ dùng \texttt{ARGS}, kiểm tra tham số thiếu, lặp hoặc lạ, và từ chối khi tệp kết
quả trùng tệp dữ liệu (kể cả khác chữ hoa thường trên Windows, đường dẫn tương đối, hard link, symlink) để
không ghi đè dữ liệu gốc.

\subsection{Kiểm thử và đối chiếu SPMF}

Mỗi bản được đối chiếu từng itemset và support với kết quả tham chiếu do SPMF sinh ra và với một chương trình
vét cạn độc lập trên dữ liệu nhỏ. Bảng~\ref{tab:fixtures} liệt kê các CSDL kiểm thử: 27 trường hợp minsup trên
7 tệp (6 CSDL khác nhau, vì hai tệp là hai bản làm sạch của cùng một dữ liệu), tổng 552 itemset tham chiếu
qua các ngưỡng; cả hai bản khớp 100\%. Ngoài ra, bộ kiểm thử đối chiếu vét cạn mọi CSDL 3 giao dịch trên 3 item ở
$\theta = 1, 2, 3$ (1.536 phép so sánh), 40 CSDL ngẫu nhiên có seed cố định, các độ dài 63/64/65 và 127/128/129
giao dịch để kiểm biên từ máy của bitset, cùng các trường hợp CSDL rỗng, ngưỡng lớn hơn $N$, item 0 và ID lớn.

\begin{table}[t]
\caption{CSDL kiểm thử và các trường hợp đối chiếu SPMF. Cột minsup: số ngưỡng đã thử; cột Itemset: tổng
số itemset tham chiếu qua các ngưỡng}
\label{tab:fixtures}
\centering
\footnotesize
\setlength{\tabcolsep}{4pt}
\begin{tabular}{@{}lrrr@{}}
\toprule
Tệp CSDL & minsup & Itemset & Khớp \\
\midrule
\texttt{dup\_items\_comments\_clean} & 3 & 30 & 100\% \\
\texttt{dup\_items\_comments\_cleaned} & 3 & 25 & 100\% \\
\texttt{dup\_tx} & 3 & 28 & 100\% \\
\texttt{rand\_dense} & 3 & 266 & 100\% \\
\texttt{rand\_sparse} & 6 & 122 & 100\% \\
\texttt{toy} & 5 & 73 & 100\% \\
\texttt{ws\_line} & 4 & 8 & 100\% \\
\bottomrule
\end{tabular}
\end{table}

% [CẦN CẬP NHẬT] Số test lấy từ lần chạy cuối cùng trước khi nộp (README cũng ghi số này).
Bộ kiểm thử chạy bằng \texttt{julia --project test/runtests.jl}; một ca kiểm thử về symlink được đánh dấu bỏ qua trên
Windows vì hệ điều hành không cho tạo symlink khi thiếu quyền.

\subsection{Độ phức tạp}

Gọi $M$ là tổng số lần xuất hiện item, $U$ số item phân biệt, $F$ số item đơn phổ biến và $w = 64$ số bit của
một từ máy. Dựng TID-list tốn $O(N + M)$ kỳ vọng; giao hai danh sách $a$, $b$ tốn $O(|a| + |b|)$. Dựng bitset
cần $O(F \lceil N/w \rceil)$ từ máy, và mỗi phép giao hoặc đếm tốn $O(\lceil N/w \rceil)$, không phụ thuộc
support. Tổng thời gian còn phụ thuộc số ứng viên được thử và độ dài kết quả; có thể có tới $2^F - 1$ itemset
phổ biến nên không thể khẳng định thuật toán tuyến tính theo $N$. Về bộ nhớ, TID-list dọc dùng $O(M + U)$, còn
bitset dọc dùng $O(F \lceil N/w \rceil)$ từ máy.

\subsection{Đo tác động của tối ưu}
\label{sec:impl-opt}

Hai CSDL tổng hợp Bernoulli với seed 42 kiểm soát mật độ: dày ($N = 12.000$, 24 item, xác suất 0,6) và thưa
($N = 20.000$, 500 item, xác suất 0,01). Thời gian gồm kiểm tra đầu vào, dựng dữ liệu dọc, duyệt sâu và tạo kết
quả, không gồm sinh dữ liệu và đọc ghi tệp; mỗi cấu hình làm nóng trước, đo 7 lần, đổi lượt chạy giữa hai bản
và báo trung vị. Bộ nhớ cấp phát là tổng cấp phát của lời gọi, khác với RAM đỉnh (peak RSS; RSS
là resident set size, bộ nhớ thường trú) của tiến trình.

\begin{table}[t]
\caption{Bản cơ bản (TID-list) và bản bitset, trung vị 7 lần đo. Mỗi ô hai số là TID-list / bitset;
cột ``Nhanh'' là số lần bitset nhanh hơn}
\label{tab:opt}
\centering
\footnotesize
\setlength{\tabcolsep}{2.5pt}
\begin{tabular}{@{}lrrrr@{}}
\toprule
Dữ liệu & Itemset & Thời gian (ms) & Nhanh & Cấp phát (MiB) \\
\midrule
Dày, 20\% & 2.324 & 394,13 / 12,54 & 31,43× & 368,96 / 21,43 \\
Dày, 35\% & 300 & 118,95 / 3,77 & 31,56× & 142,20 / 3,82 \\
Thưa, 0,5\% & 500 & 204,17 / 117,82 & 1,73× & 17,48 / 317,37 \\
Thưa, 1\% & 248 & 53,81 / 20,55 & 2,62× & 7,17 / 78,30 \\
\bottomrule
\end{tabular}
\end{table}

Bảng~\ref{tab:opt} cho thấy trên dữ liệu dày, bitset vừa nhanh hơn khoảng 31 lần vừa cấp phát ít hơn. Trên dữ
liệu thưa, bitset vẫn nhanh hơn nhưng cấp phát nhiều hơn hẳn (317,37 so với 17,48 MiB ở minsup 0,5\%). RAM đỉnh
của hai bản gần bằng nhau (khoảng 292 đến 325 MiB) vì phần lớn là của tiến trình Julia. Mục~\ref{sec:eval} kiểm
tra hai biểu diễn này ở quy mô lớn hơn trên dữ liệu thật.
